%% Preface -------------------------------------------------------------------
\chapter{Preface}

The \appname{} began as a tool for looking at the solar radio dynamic spectra recorded by
the e-CALLISTO network and for measuring the drift of Type~II bursts. Over successive
releases it has grown into a framework that also retrieves archive data from several
radio instruments, places a radio event in its wider solar and geophysical context,
analyzes solar images from space-borne observatories, and reconstructs coronal mass
ejections in three dimensions. This guide describes all of that functionality as it
exists in version~\appversion.

\section*{About this guide}

This is a \emph{user guide}: a reference to what each window, menu, panel and control does,
what it expects as input, what it produces, and which limits apply. It is not a tutorial
and does not teach solar radio physics or image analysis. Where an operation consists of
several steps, the steps are given as a short numbered procedure; otherwise each feature
is described in the place where you meet it in the interface. Short \emph{Background}
boxes summarize the models behind computed quantities so that results can be interpreted
and reported correctly; the full set of equations is collected in
Appendix~\ref{app:formulas}.

The guide is written for observers, students and researchers who work with e-CALLISTO
data and related space-weather observations. No programming knowledge is required.

\section*{How this guide is organized}

\begin{description}
  \item[Part I, Getting Started] introduces the software, explains installation on
    Windows, macOS and Linux, tours the main window, and describes the help, update,
    citation and bug-report tools.
  \item[Part II, Working with Dynamic Spectra] covers opening and combining FITS files, the
    sidebar processing panels, display, RFI cleaning, annotation and measurement tools,
    the Type~II drift and shock analysis, band-splitting analysis, projects and export.
  \item[Part III, Radio Data Access] describes the e-CALLISTO FITS Downloader, the
    Multi-Station Comparison workspace, and the Learmonth and STEREO/SWAVES loaders.
  \item[Part IV, Solar-Event Context] describes the GOES X-ray, GOES proton, Dst, Kp and
    SOHO/LASCO CME catalog viewers, and the SunPy Multi-Mission Explorer.
  \item[Part V, Solar Image Analysis] covers the multi-mission imaging workspace: data
    retrieval, display and processing, measurement and CME kinematics, and potential-field
    source-surface modeling.
  \item[Part VI, Three-Dimensional CME Reconstruction] describes the GCS CME Fitting window
    for flux-rope and shock reconstruction from several coronagraph viewpoints.
  \item[The appendices] list keyboard shortcuts, file formats and outputs, data sources,
    models and formulas, and troubleshooting advice, followed by a glossary, the
    bibliography and an index.
\end{description}

\section*{Conventions}

Table~\ref{tab:conventions} lists the typographic conventions used throughout the guide.

\begin{table}[!htbp]
  \centering
  \caption{Typographic conventions used in this guide.}
  \label{tab:conventions}
  \small
  \begin{tabularx}{\linewidth}{@{}P{0.34\linewidth}L@{}}
    \toprule
    \thead{Example} & \thead{Meaning} \\
    \midrule
    \menu{File > Export As > Export Figure} & A command reached through a menu and its submenus. \\
    \gui{Subtract Background} & A button, field, check box, tab, panel or window as labeled in the interface. \\
    \kbd{Ctrl+S} & A keyboard shortcut. On macOS, \kbd{Ctrl} in a shortcut corresponds to the Command key \kbd{Cmd}. \\
    \file{project.efaproj} & A file name, file extension or folder. \\
    \callout{3} & A numbered callout in a screenshot, explained in the accompanying table. \\
    \textsf{\textbf{\textcolor{accent}{Note}}}, \textsf{\textbf{\textcolor{tipc}{Tip}}}, \textsf{\textbf{\textcolor{cautionc}{Caution}}} & Boxes with additional detail, a quicker way of doing something, or a warning about lost work or results that need care. \textsf{\textbf{\textcolor{scic}{Background}}} boxes summarize the model behind a computed quantity. \\
    \bottomrule
  \end{tabularx}
\end{table}

The screenshots were taken with the macOS version of the software in the light theme.
On Windows and Linux the window frames, fonts and menu placement follow the conventions
of the operating system, and in the dark theme the colors differ, but the controls and
their behavior are identical on all platforms.

\section*{Acknowledgments}

The \appname{} relies on data from the e-CALLISTO network and its archive, and on the many
public archives listed in Appendix~\ref{app:data-sources}. It is built on open-source
scientific software, including Qt for Python (PySide6), NumPy, SciPy, Matplotlib,
PyQtGraph, Astropy \parencite{astropy2022}, SunPy \parencite{sunpy2020}, aiapy
\parencite{barnes2020}, sunkit-magex \parencite{stansby2020} and reproject. Their authors and
maintainers are gratefully acknowledged.

\begin{flushright}
  \textit{Sahan S Liyanage}\\
  \textit{Colombo, October 2026}
\end{flushright}
